\section{Introduction}
\label{sec:introduction}

Vision-language-action models (VLAs) achieve impressive robotic manipulation performance by fine-tuning large language models on robot demonstration data~\cite{brohan2023_rt2,kim2024_openvla}.
The dominant architecture grafts a frozen vision encoder onto a pretrained LLM backbone, then fine-tunes the system end-to-end to map visual observations and language instructions to continuous actions~\cite{karamcheti2024}.
An implicit assumption underlies this paradigm: fine-tuning injects action-relevant representational knowledge into the LLM.
But where do action representations actually come from?
Do they originate in the vision encoder's pretrained features, emerge through LLM processing, or arise from task-specific fine-tuning?

Recent work reveals cracks in this assumption.
Grant et al.~\cite{grant2026} showed that the visual pathway dominates action generation across multiple VLA architectures, with ablations to the language stream producing minimal behavioral degradation.
Independently, Hou et al.~\cite{hou2026} and Fei et al.~\cite{fei2025} demonstrated that VLAs largely ignore language instructions during policy execution.
These findings challenge the prevailing narrative but remain at the behavioral level---they reveal \emph{what} drives VLA outputs without explaining \emph{where} action-relevant structure resides in the representation space.
No prior work has performed a representational decomposition that measures each component's independent contribution to action-relevant structure.
Without this, we cannot determine whether fine-tuning truly injects action knowledge or merely learns to route pre-existing structure.

We present the first systematic representational and functional analysis of a VLA.
Taking OpenVLA-7B~\cite{kim2024_openvla} as a case study, we decompose its processing pipeline into constituent components: the frozen SigLIP vision encoder~\cite{zhai2023_siglip} with MLP projector, every layer of the Llama-2 7B backbone, and the readout token position.
We apply Representational Similarity Analysis (RSA)~\cite{kriegeskorte2008} to measure geometric alignment with the action space and train functional probes to measure action decodability.
A $2 \times 2$ factorial design crossing training status (trained vs.\ untrained LLM) with instruction presence (with vs.\ without) isolates each factor's contribution.
External baselines---DINOv2~\cite{oquab2024_dinov2} (vision-only) and CLIP ViT-L/14~\cite{radford2021_clip} (vision-language pretrained, no VLA training)---anchor our comparisons.

Three findings emerge from experiments on the LIBERO-Goal benchmark~\cite{liu2023_libero}.
\textbf{First}, frozen vision-language encoders provide the dominant action-relevant geometric structure.
SigLIP+Projector achieves the highest action-space alignment (RSA $\rho = 0.615$), followed by CLIP ($\rho = 0.557$), both far exceeding the vision-only DINOv2 ($\rho = 0.388$).
Yet DINOv2 matches SigLIP+Projector in action prediction capacity ($R^2 = 0.712$ vs.\ $0.706$, $p = 0.22$), revealing a geometric-functional dissociation: VL pretraining provides geometric alignment with the action space without increasing raw decodability.
\textbf{Second}, the LLM degrades distributed geometric alignment (RSA $0.615 \to 0.501$) but concentrates a functionally superior signal at the readout token ($R^2 = 0.731$, exceeding the projector's $0.706$).
This concentration is training-dependent: an untrained LLM's readout achieves only $R^2 = 0.686$, falling below the projector baseline.
\textbf{Third}, training teaches routing, not creating.
Distributed RSA is training-invariant ($\Delta = 0.021$, 95\% CI includes zero, $p = 0.20$), while readout concentration shows a large training effect ($\Delta = 0.045$, $p < 0.0001$).
Language instructions affect neither metric.

These results reframe how we understand VLA representations and carry direct design implications: vision encoder pretraining quality---not LLM fine-tuning---is the primary lever for action-relevant structure~\cite{haon2025}.
Our contributions are:
\begin{enumerate}[leftmargin=*,itemsep=2pt,topsep=2pt]
    \item The first systematic representational decomposition of a VLA, quantifying each component's contribution to action-relevant structure with both geometric (RSA) and functional (probe $R^2$) metrics.
    \item Discovery of a geometric-functional dissociation: vision-language pretraining provides geometric alignment with the action space beyond what raw information content predicts, while vision-only features are functionally rich but geometrically dissimilar.
    \item Evidence that LLM fine-tuning teaches information routing---concentrating pre-existing vision-encoder structure at the readout position---rather than creating new action-relevant representations.
\end{enumerate}
